% !TeX program = pdflatex
% !TeX encoding = UTF-8
% P3：真空微重力环境下液桥诱导风化层增强的锚定可行性与失效边界。
% 当前状态：研究规划；液体、颗粒、装置和可行性均未确定。
% 本篇不以最终证明可用为前提；若能量化不可行边界与控制机制，否定结果有效。
% 不能把“湿颗粒强度增加”直接等同“飞行器锚定成功”，也不能将注入固结整体
% 宣称为首次；需与泡沫、胶黏/固化和既有湿颗粒研究区分。
%
% 候选假设：
% H1 低挥发、良好润湿液体可形成持久液桥并提高低围压颗粒表观黏聚。
% H2 含液量存在最优区间；过少连接不足，过多导致饱和、迁移或润滑。
% H3 局部强度提高后，整体失效转移到锚具接口或增强区外部，承载出现上限。
% H4 微重力提高毛细力相对重要性，但不自动扩大外围失效面积。
% H0 在现实剂量、等待时间、环境与注入扰动下没有系统级净锚定收益。
%
% 图路线：F01 载荷链；F02 液体保持/浸润；F03 局部增强；F04 整体承载；
% F05 失效面；F06 标度；F07 环境/资源可行域。

\documentclass[11pt]{article}
\usepackage[T1]{fontenc}
\usepackage[utf8]{inputenc}
\usepackage{amsmath,amssymb}
\usepackage{booktabs}
\usepackage{graphicx}
\usepackage[hidelinks]{hyperref}
\usepackage[margin=1in]{geometry}

\title{Feasibility and Failure Limits of Capillary-Bridge-Induced\\
Regolith Anchoring under Vacuum and Microgravity}
\author{Authors to be confirmed}
\date{Planning manuscript --- evidence pending}

\newcommand{\planned}[1]{\par\noindent\textit{Planned work: #1}\par}
\newcommand{\figureplaceholder}[1]{%
  \fbox{\parbox{0.91\linewidth}{\centering
  \vspace{0.018\textheight}\textbf{Planned evidence --- no results yet}\par
  \smallskip #1\par\vspace{0.018\textheight}}}}
\newif\ifBibliographyReady
\BibliographyReadyfalse

\begin{document}
\maketitle

\begin{abstract}
This planning manuscript asks whether capillary cohesion generated by a
low-volatility liquid can be converted from local granular strengthening into
system-level anchoring on a rubble-pile asteroid. The proposed study follows the
complete load path from liquid retention and wetting through particle bridges,
anchor coupling, reinforced-zone failure, and detachment from untreated regolith.
It will map useful and infeasible regimes under environmental and resource
constraints. No candidate liquid or anchoring benefit has yet been demonstrated.
\end{abstract}

\section{Scientific question and prior-art boundary}
% 精确缺口：毛细局部增强能否产生系统级锚力，以及真空/温度/时间和外围失效边界。
% 比较湿颗粒、非饱和土、液体桥、泡沫/胶黏/固化锚定、颗粒注入增强。
% 若实际机理变为化学固化/冻结/胶黏，必须重新命名，不继续称为液桥。
\planned{Verify prior work and define the difference between reversible capillary
bridges, frozen or cured bonds, foam stabilization, adhesive anchoring, and
pressure-driven injection. State a falsifiable system-level question.}

\section{Load path and failure hierarchy}
% 载荷链：锚具→锚具/增强体界面→增强体内部→增强体/原风化层外边界。
% 整体承载受最弱环节限制；研究拉、剪、矩，不能只测湿料抗压。
A first-order upper bound is
\begin{equation}
F_{\mathrm{anchor}}\leq
\min\!\left(F_{\mathrm{interface}},F_{\mathrm{reinforced}},
F_{\mathrm{outer}}\right),
\end{equation}
where the three terms describe anchor--zone coupling, internal failure of the
reinforced zone, and detachment through surrounding untreated regolith.

\planned{Define the full force and moment path, observable failure surfaces, and
conditions under which strengthening simply moves failure to another interface.}

\begin{figure}[htbp]
\centering
\figureplaceholder{P3-F01: Liquid delivery, capillary network, anchor coupling,
reinforced volume, surrounding regolith, external load, and competing failures.}
\caption{Planned system boundary and load path.}
\end{figure}

\section{Liquid retention, wetting, and transport}
% 变量：蒸气压、表面张力、接触角、黏度、温度、质量损失、孔隙、剂量、注入速率。
% 真空不消除表面张力，但改变液体保持；普通喷注可能吹散低强度颗粒。
% 记录质量/形貌随时间，区分初始注入和毛细自发浸润。
\planned{Characterize candidate liquid--mineral pairs over the required pressure,
temperature, and time. Measure contact angle, surface tension, viscosity, mass
loss, infiltration morphology, and injection disturbance rather than assuming
nominal properties persist in the test environment.}

\begin{figure}[htbp]
\centering
\figureplaceholder{P3-F02: Retained liquid mass, wetting front, spatial saturation,
and characteristic times against pressure, temperature, dosage, and grain state.}
\caption{Planned environmental persistence and transport evidence.}
\end{figure}

\section{Particle-scale and material-scale capillary cohesion}
% 量级 c_cap ~ C gamma cos(theta)/d；C 需由级配、孔隙、桥数、饱和度标定。
% 直接/独立液桥证据：成像、单桥力、液体分布或可验证的本构响应。
% 对照注入扰动、压实、重排和润滑，不能仅干/湿比较。
Use the scaling
\begin{equation}
c_{\mathrm{cap}}\sim C\frac{\gamma\cos\theta}{d}
\end{equation}
only as a hypothesis to test, where $C$ must account for grading, porosity,
liquid content, contact network, and the relevant loading path.

\planned{Measure tensile, shear, and relevant cyclic response at low normal
stress. Test liquid-content optima, grain size and shape, hysteresis, rate, and
aging while obtaining direct or independent evidence for the capillary mechanism.}

\begin{figure}[htbp]
\centering
\figureplaceholder{P3-F03: Local strength and deformation against capillary
scales, liquid content, grain properties, normal stress, time, and controls.}
\caption{Planned local strengthening and constitutive evidence.}
\end{figure}

\section{System-scale anchoring tests}
% 同一锚具/床制备下测试拉拔、剪切、力矩及循环。记录完整力位移、机体/床运动、
% 液体量、等待时间和所有失败。局部材料试样与整体锚体必须分开。
% 最小对照：原始干床；匹配针具进入/注入机械扰动但无液体；注液；必要时低润湿
% 或相同压实度。只改液体时多物性同时变化，不能把所有差异归表面张力。
\planned{Compare complete anchors in independently prepared beds using dry,
matched-disturbance, and liquid-treated conditions. Measure pull, shear, moment,
and cyclic behavior together with reinforced-zone geometry and resource cost.}

\begin{figure}[htbp]
\centering
\figureplaceholder{P3-F04: Local material gain versus complete anchor capacity,
including conservative quantiles, dosage, waiting time, injection reaction,
energy, and all attempted deployments.}
\caption{Planned test of whether local strengthening becomes system-level gain.}
\end{figure}

\section{Failure localization and mechanism tests}
% 通过透明/成像/标记/挖掘后观察、数字图像/断层等可用方法定位失效。
% 干预增强区半径/深度/分支/锚接口，检查外围失效假设。不能仅凭裂后照片断言因果。
\planned{Locate initiation and propagation of failure, align diagnostics in time,
and use controlled changes in reinforced geometry or interface coupling to test
whether the outer untreated boundary controls capacity.}

\begin{figure}[htbp]
\centering
\figureplaceholder{P3-F05: Internal, anchor-interface, and outer-boundary failure
modes with controlled interventions and ensemble statistics.}
\caption{Planned causal evidence for the limiting load-transfer stage.}
\end{figure}

\section{Scaling and microgravity transfer}
% 量纲候选：Bond 数、毛细数、粒径/增强尺度、含液量、黏聚/任务应力、时间比等。
% 地面真空不等于微重力；机器人卸载不减少颗粒自重。分开讨论毛细/重力相对尺度、
% 低围压、制样历史和实际低重力验证。避免将固定 g 扫描当全部相似律。
\planned{Derive and test dimensionless controls that distinguish capillary,
gravity, viscous, inertial, geometric, and time scales. Quantify what terrestrial
vacuum tests validate and what remains an extrapolation to reduced gravity.}

\begin{figure}[htbp]
\centering
\figureplaceholder{P3-F06: Scaling collapse or documented failure across grain,
reinforced-zone, gravity/normal-stress, rate, and liquid-property conditions.}
\caption{Planned physical relation and transfer limits.}
\end{figure}

\section{Feasibility, resources, and negative outcomes}
% 用 P0 任务场景和 P1 路径约束，综合整体承载低分位、液体质量、等待、注入反力、
% 热控/储存、污染、持续时间。零净收益或不可行边界可成篇，但需机制+定量上限。
\planned{Map beneficial, infeasible, and unresolved regimes under mission load,
liquid mass, deployment time, injection disturbance, storage, thermal, and
retention constraints. Compare with an appropriate mechanical or solidification
baseline without assuming capillary anchoring is superior.}

\begin{figure}[htbp]
\centering
\figureplaceholder{P3-F07: Mission-conditioned regime map showing net benefit,
outer-boundary limits, environmental loss, excessive resource cost, and untested
regions with uncertainty.}
\caption{Planned feasibility and no-go boundaries.}
\end{figure}

\section{Completion and stop conditions}
% 成篇：机理有独立证据；局部到整体载荷链；环境/剂量/时间；失效边界；资源评价。
% 即使整体无增益，只要确定为何、何处、上限及适用范围，仍可成篇。
% 止损系统开发：只在空气有效；仅改善均值；剂量/等待不可接受；机理实为固化。
\planned{Complete the paper when the capillary mechanism, load-transfer chain,
and feasibility boundary are quantitatively supported, whether the resulting
mission verdict is positive or negative. Stop capillary-system development when
the mechanism, persistence, or resource case fails, and rename the route if the
actual support comes from curing, freezing, or adhesion.}

\section*{Data and software availability}
\planned{Archive liquid characterization, environmental histories, sample and
bed preparation, raw load records, images, failures, models, and analysis code.}

\ifBibliographyReady
  \bibliographystyle{plain}
  \bibliography{references}
\else
  \section*{References to verify}
  Verified sources on wet granular mechanics, vacuum-compatible liquids,
  asteroid regolith, injection stabilization, and anchoring will be added after
  the mechanism and comparison boundary are frozen.
\fi

\end{document}
